\subsection{F(h), added item (2): the punctured completion and the end}

\paragraph{F.40 (Proved). Nonarchimedean invariance, with the missing hypothesis.}
Let $k$ be any field, equipped with the trivial absolute value, let $X$ be
a separated $k$-scheme of finite type, and let $U\subset X$ be a dense open
subscheme with complement $D$. Denote by $Y^{\beth}$ the generic fibre of
the scheme $Y$ regarded as a formal $k$-scheme with the discrete topology.
The analytic punctured tube is
\[
 W(X,D)=X^{\beth}\setminus\bigl(U^{\beth}\cup D^{\beth}\bigr).
\]
It is the punctured generic fibre associated with $\widehat X_D$.
For every proper morphism $X'\to X$ which is the identity on $U$, the
induced morphism $W(X',X'\setminus U)\to W(X,D)$ is an isomorphism of
$k$-analytic spaces. If $X$ is proper over $k$, then
\[
                 W(X,D)=U^{\mathrm{an}}\setminus U^{\beth}=:U_\infty^{\mathrm{na}}.
\]
Consequently it is independent of the proper compactification, without
smoothness, normal-crossings, or characteristic-zero assumptions. The
phrase ``with no hypothesis'' is false if it also removes properness.

\begin{proof}
A point of $Y^{\beth}$ is a multiplicative seminorm whose valuation ring
has a centre on $Y$; separatedness makes the centre unique. A point of
$W(X,D)$ has its generic point in $U$ and its centre in $D$. The valuative
criterion for a proper $X'\to X$ gives a unique lifted centre, necessarily
in $X'\setminus U$. Thus both tubes are the same analytic domain of
$U^{\mathrm{an}}$, with the same induced analytic structure. If $X$ is
proper, every analytic point has a centre, so $X^{\beth}=X^{\mathrm{an}}$
and $D^{\beth}=D^{\mathrm{an}}$, giving the formula. This is also the
precise statement of Ben-Bassat--Temkin, \S3.3.2 and Proposition~4.5.1
\cite{BenBassatTemkin2012}.

For necessity, take $U=\mathbb G_m$ over an algebraically closed field.
Its embedding in $\mathbb A^1$ has boundary $\{0\}$ and tube
$0<|t|<1$, a single open interval. Its proper compactification
$\mathbb P^1$ has two boundary points and tube
$\{0<|t|<1\}\amalg\{1<|t|<\infty\}$, two open intervals. These are not
even equally connected. The first ambient space is a partial
compactification, so it is not a counterexample to the properly stated
theorem.

The notation $\widehat X_D\setminus D$ on MPL~149, fol.~66, scan~133
must be given this punctured-generic-fibre meaning: literally deleting
$D$ from the underlying topological space of the formal completion gives
the empty space. In particular, the invariant does not assert that the
formal schemes $\widehat X_D$ themselves are isomorphic.
\end{proof}

\paragraph{F.41 (Refuted). Stronger replacement: the archimedean link is invariant too.}
The proposed failure of compactification invariance for the archimedean
punctured neighbourhood when $X$ is compact is false. For every compact
Hausdorff compactification $U\subset X$, with $D=X\setminus U$, its
punctured-neighbourhood pro-space is canonically
\[
 \{N\setminus D\}_{N\supset D}
       =\{U\setminus K\}_{K\subset U\ \mathrm{compact}}.
\]
Here $N$ ranges over open neighbourhoods, and the equality is an
isomorphism of indexing systems, not just equality of cohomology.
For a smooth proper complex variety with a simple normal-crossings
boundary, this pro-homotopy type is constant, represented by the boundary
$L$ of its real oriented blow-up, after rounding corners. Thus the
homotopy type of $L$ is independent of that compactification. In general
the natural map from this end to $U$ is not a homotopy equivalence.

\begin{proof}
If $N$ is open and contains $D$, then $K=X\setminus N$ is a compact
subset of $U$ and $N\setminus D=U\setminus K$. Conversely, every compact
$K\subset U$ is closed in the Hausdorff space $X$, so $N=X\setminus K$
is such a neighbourhood. These assignments are inverse and preserve the
refinement maps. The same proof applies to punctured analytic
neighbourhoods, since these are open subsets of the common complex
manifold $U$.

In a normal-crossings coordinate chart $z_1\cdots z_r=0$, write
$z_j=\rho_j e^{i\theta_j}$. The real oriented blow-up replaces the
vanishing coordinates by $\rho_j\geq0$ and their angles. Its local model
is $[0,\infty)^r\times(S^1)^r\times\mathbb R^{2n-2r}$, so it is a
compact manifold with corners with interior $U$. Rounding corners and
taking a boundary collar gives a cofinal family of end neighbourhoods
$L\times(0,\epsilon)$. All refinement inclusions are homotopy
equivalences, which proves the constant-model assertion. The constant
embedding of the homotopy category into its pro-category is fully
faithful; hence two such models have equivalent homotopy types. Weber's
\S2 identifies the corresponding cohomology directly with
$\varinjlim_K H^*(U\setminus K)$ \cite{Weber2012}; the argument here
retains the whole pro-homotopy type. The explicit failures of $L\to U$
are proved in F.42.
\end{proof}

\paragraph{F.42 (Proved). Both requested examples and the three different invariants.}
For $(X,D)=(\mathbb P^1,\{0,1,\infty\})$ over $\mathbb C$, one has
\[
 L\simeq S^1\amalg S^1\amalg S^1,\qquad
 U(\mathbb C)\simeq S^1\vee S^1,
\]
while over trivially valued $\mathbb C$ the underlying topological space
of $U_\infty^{\mathrm{na}}$ is three disjoint open intervals. For
$(X,D)=(\overline M_{0,5},\partial\overline M_{0,5})$ the invariants are
\[
\begin{array}{c|rrrr}
 &b_0&b_1&b_2&b_3\\\hline
 U(\mathbb C)&1&5&6&0\\
 L&1&11&11&1\\
 U_\infty^{\mathrm{na}}\text{ (underlying topology)}&1&6&0&0.
\end{array}
\]
All displayed integral homology groups are free. Both end constructions
are independent of a proper compactification in the precise senses of
F.40--F.41; neither the archimedean link nor the underlying topological
space of the nonarchimedean tube is the homotopy type of $U(\mathbb C)$.
The last row concerns Berkovich topology, not its etale topos.

\begin{proof}
For the three-punctured sphere, small punctured disks retract to three
circles and a pair of pants retracts to a wedge of two circles. A
multiplicative seminorm in the formal punctured disk satisfies
$|\sum_{n\geq m}a_nt^n|=r^m$, where $a_m\ne0$ and $0<r<1$; a unit and
its inverse are bounded, so they have norm one. Thus $r$ parametrizes an
open interval. The three boundary points give the asserted disjoint
union. Its normalized quotient by rescaling valuations is three points.

Normalize three marked points of a smooth five-pointed curve to
$0,1,\infty$. Then
\[
 U(\mathbb C)=\{(x,y)\in(\mathbb C\setminus\{0,1\})^2:x\ne y\}.
\]
Projection to $x$ is a locally trivial bundle, with base the plane minus
two points and fibre the plane minus three points. A local trivialization
is obtained by an isotopy moving $x$ in a small disk disjoint from $0,1$.
Both base and fibre retract to graphs, so the homotopy exact sequence
proves that $U$ is a $K(\pi,1)$. Monodromy preserves each oriented
puncture-meridian in first homology. The integral Serre spectral sequence
therefore has only its $0$ and $1$ rows and columns, all free, with ranks
given by $(1+2t)(1+3t)$. It has no possible differential, yielding the
first row of the table.

The boundary divisors are indexed by two-element subsets of
$\{1,\ldots,5\}$. Two distinct divisors meet exactly when these subsets
are disjoint, as follows by drawing the unique stable three-component
tree with the two prescribed nodes. They intersect transversely once;
there are no triple intersections. Thus the incidence graph is the
Petersen graph, with $10$ vertices and $15$ edges, and its first Betti
number is $15-10+1=6$. Every component is a projective line of
self-intersection $-1$. These geometric descriptions also follow from
the identification with the blow-up of $\mathbb P^2$ in four points,
where the boundary consists of four exceptional curves and the six
transforms of their joining lines; see Dolgachev--Farb--Looijenga,
\S\S2.1--2.2 \cite{DolgachevFarbLooijenga2018}.

For the nonarchimedean row, the normal-crossings skeleton theorem gives
a deformation retraction of the tube onto
$|\Delta(D)|\times\mathbb R_{>0}$: in a chart its coordinates are the
nonnegative numbers $-\log|z_j|$, not all zero, and normalization divides
by their sum. The coordinate retractions glue under etale changes of
normal-crossings charts. This is Thuillier's Theorem~3.25 and
Proposition~4.5 in the opened 2006 preprint \cite{Thuillier2006}.
Consequently the tube has the homotopy type of the Petersen graph.

For the archimedean row, let $N$ be a closed regular neighbourhood of
$D$ in $X$. It is obtained by plumbing the disk normal bundles of the
ten components, and retracts to $D$. It follows that
$H_1(N;\mathbb Z)=\mathbb Z^6$, $H_2(N;\mathbb Z)=\mathbb Z^{10}$,
and $H_3(N;\mathbb Z)=0$. The boundary $L=\partial N$ is a connected
closed oriented three-manifold. Let $Q$ be the intersection matrix of
the ten components. Its diagonal is $-1$ and its off-diagonal entries
are the adjacency matrix of the Petersen graph. Poincare--Lefschetz
duality and the homology exact sequence of $(N,L)$ give
\[
 0\longrightarrow\operatorname{coker}Q\longrightarrow H_1(L;\mathbb Z)
       \longrightarrow\mathbb Z^6\longrightarrow0.
\]
To compute the cokernel integrally, use the basis $H,E_1,\ldots,E_4$ of
$H_2(X;\mathbb Z)$ with form $\operatorname{diag}(1,-1,-1,-1,-1)$.
The boundary classes are $E_i$ and $H-E_i-E_j$. They generate this
unimodular lattice, so their column map $V:\mathbb Z^{10}\to\mathbb Z^5$
is onto. Since $Q=V^t\operatorname{diag}(1,-1,-1,-1,-1)V$, its image is
the primitive rank-five sublattice $\operatorname{im}V^t$. Hence
$\operatorname{coker}Q=\mathbb Z^5$ and $H_1(L)=\mathbb Z^{11}$.
Duality now gives $H_2(L)=\mathbb Z^{11}$ and $H_3(L)=\mathbb Z$.
The attached program independently computes the graph incidence rank
and Smith normal form of $Q$.

Already connectedness rules out $L\simeq U$ in the curve example; in
the surface example the nonzero fundamental class in $H_3(L)$ rules it
out. Neither fact contradicts compactification invariance. The
manuscript's cautious assertion on scan~133 is therefore supported
when ``punctured completion'' denotes the nonarchimedean tube or the
archimedean punctured germ; its identification with the open variety is
not asserted by that sentence and is false in these examples.
\end{proof}
